\section{Conclusion}\label{sec:conclusion}
Intensity interferometry (II) is re-emerging as a promising technique for optical long-baseline interferometry. However, unlike conventional amplitude interferometry, II measures the correlation of photon arrival times and provides information on the squared visibility, resulting in the loss of phase information. Recovering the source morphology from such measurements therefore remains an important challenge.

In this work, we have explored the feasibility of using a conditional Generative Adversarial Network (cGAN), based on the Pix2Pix framework, for image reconstruction from II data. The model was trained using paired synthetic stellar images and their corresponding II data, with the ground-truth sky images serving as targets and the II measurements as conditional inputs. The generator learns to reconstruct images consistent with the input II data, while the discriminator encourages the generation of realistic stellar intensity distributions. The reconstruction performance was evaluated by comparing the image moments of the generated and ground-truth images, including the monopole and second- and third-order moments. The results demonstrate that the proposed approach can recover the overall morphology and brightness distribution of the synthetic rapidly rotating stars considered in this study.

{\bf The present work provides a controlled proof of concept rather than a comprehensive assessment of image reconstruction for realistic stellar surfaces. The adopted synthetic images have smooth brightness distributions and do not incorporate features such as starspots, granulation, realistic limb darkening, or tidal deformation in interacting binary systems. Moreover, the adopted brightness prescription is phenomenological and should not be interpreted as a complete physical model of gravity darkening. Thus, the demonstrated reconstruction capability is limited to the synthetic benchmark investigated here.

Several extensions are required to establish the applicability of this approach to more realistic astrophysical sources. Future work should incorporate physically motivated gravity-darkening prescriptions, wavelength-dependent limb darkening, surface inhomogeneities, and tidal distortions. It will also be important to investigate the effects of Fourier-plane coverage, telescope-array geometry, observing time, and detector efficiency on reconstruction performance. In particular, the improved spatial-frequency sampling potentially offered by larger IACT arrays and future facilities such as CTAO may provide valuable opportunities to investigate more complex stellar intensity distributions. The effects of these factors will need to be assessed through systematic simulations.

Although the relatively simple sources considered here could also be investigated using conventional parameter-fitting techniques, the objective of this study was to explore machine-learning-based image reconstruction from sparsely sampled II data. The present results do not establish the superiority of cGAN over conventional parameter-fitting methods, nor do they quantify the specific contribution of adversarial learning relative to an L1-only U-Net. Comparisons with these approaches, along with investigations of alternative network architectures and loss functions, will be necessary to assess their relative advantages and limitations.}

Overall, this study represents an initial step toward applying machine-learning techniques to image reconstruction in II. Further development and validation using more realistic stellar models, improved noise descriptions, and different telescope-array configurations will help determine the potential of this approach for studying the morphology and surface-brightness distributions of rapidly rotating and interacting stars.